\chapter{Conclusion}
\label{ch:concl}

\section{Summary}
\label{sec:summary}

This thesis has addressed the question of how manipulation demonstrations
recorded passively, with a human using their own hands and no teleoperation
interface, can be converted into datasets suitable for training visuomotor
imitation policies. The work has argued that the two established approaches to
this problem are each incomplete: online teleoperation systems have converged
on a sound retargeting formulation---a single weighted cost functional
combining data fidelity, temporal smoothness and kinematic constraints---but
produce no dataset, while offline video-to-robot pipelines produce datasets but
retarget either not at all, through sparse object keyframes, or through a
sequential filter--solve--filter pattern whose smoothing is applied in the wrong
space to guarantee the property it is intended to provide.

\viki\ has been presented as a system occupying the intersection. It captures
demonstrations with a metrically calibrated multi-camera RGB-D rig anchored to
a ChArUco workspace frame; represents the hand minimally, as a wrist pose in
$\SE$ with a binary gripper state; stores each demonstration both relative to
the manipulated object and relative to the workspace anchor; retargets by
minimising one weighted functional in which the data term is scaled by
per-frame perception confidence and robustified by a Huber loss, and in which
velocity and acceleration penalties act directly in joint space; validates the
result by executing it on the physical manipulator before export; and writes a
LeRobot-compatible dataset containing the proprioception the robot actually
attained rather than the states the optimiser predicted.

\section{Answers to the Research Questions}
\label{sec:answers}

\paragraph{RQ1.} \ph{[State the measured effect of the single-functional
formulation on fidelity and smoothness relative to the sequential baseline, and
whether the trade-off frontier moved rather than merely shifted along.]}

\paragraph{RQ2.} \ph{[State the achieved wrist-pose accuracy and how the
residual decomposes across landmark noise, calibration error and
representation dependency, per Table~\ref{tab:errbudget}.]}

\paragraph{RQ3.} \ph{[State the proportion of episodes rejected by replay
screening, the dominant rejection causes, and the resulting change in the
executable fraction of the exported dataset.]}

\section{Contributions}
\label{sec:contrib-final}

The contributions of this work are the transfer of the single-cost-functional
retargeting formulation from the online teleoperation regime, for which it was
developed, into the offline object-centric dataset-generation regime, where the
review of Chapter~\ref{ch:lit} found it unused; the coupling of that
formulation to a metrically calibrated multi-view RGB-D front end with
confidence-weighted fusion and explicit cross-device timestamp handling; the
use of physical replay as a pre-export dataset-curation stage rather than a
post-training evaluation; and an error decomposition that separates the three
independent contributions to residual pose error rather than reporting an
aggregate.

It is worth restating what is \emph{not} claimed. The cost functional is not
novel as an operator; its components appear in \cite{anyteleop, dexflow,
opentv, dexteleop0}, and a claim of algorithmic novelty would be refuted by
that prior art. Object-centric representation, ChArUco anchoring,
confidence-weighted depth fusion and host-clock synchronisation are all
established practice. The contribution is the conjunction and the regime
transfer, which the positioning analysis of Section~\ref{sec:positioning} found
unoccupied.

\section{Limitations}
\label{sec:concl-limits}

The system is restricted to parallel-jaw targets, rigid objects, single-arm
quasi-static tabletop tasks and a fixed workspace; it performs kinematic
transfer only, estimating no forces; its accuracy is bounded below by
centimetre-level cross-camera agreement; and its evaluation uses a single
demonstrator on a constrained task class. These are stated in full in
Section~\ref{sec:limitations}.

\section{Future Work}
\label{sec:future}

Several directions follow from the limitations.

\begin{enumerate}[leftmargin=2em]
  \item \emph{Dexterous targets.} Extending beyond a binary gripper requires
        reinstating a finger-level representation. The ablation evidence of
        \cite{keyobj} suggests that the wrist objective should remain weakly
        weighted relative to fingertip objectives, which would change the
        balance of terms in \eqref{eq:cost} rather than its structure.
  \item \emph{Deformable objects.} Replacing the rigid object pose of
        \eqref{eq:objrel} with a flow-based object representation
        \cite{novaflow} would remove the rigidity assumption at the cost of a
        less compact transfer rule.
  \item \emph{Closing the replay loop.} Currently, replay failures trigger
        bounded re-solution with hand-adjusted weights. Learning the weight
        adjustment from the failure mode would make the loop automatic.
  \item \emph{Cross-embodiment transfer.} The object-relative representation is
        embodiment-agnostic by construction; the same recorded demonstrations
        could be retargeted to the KUKA iiwa14 available in the laboratory, and
        the degree to which policy performance transfers would test the
        representational claim directly.
  \item \emph{Scaling demonstrator diversity.} Recording from multiple
        demonstrators with differing hand morphologies would test whether the
        fixed hand-to-end-effector offset $\bT^{H}_{E}$ of
        \eqref{eq:transfer} requires per-demonstrator calibration.
  \item \emph{Downstream curation.} The confidence weights and replay residuals
        retained in the exported dataset support curation experiments in the
        manner of \cite{demoscore} without re-recording; whether the replay
        residual is a better curation signal than a learned rollout classifier
        is an open question.
\end{enumerate}

\section{Concluding Remarks}
\label{sec:remarks}

The practical claim of this work is modest and, for that reason, defensible:
a formulation already accepted in one part of the robot-learning literature
works in another part where it had not been applied, and the infrastructure
required to apply it---metric multi-view sensing, honest cross-device
synchronisation, physical validation before export---is buildable at the scale
of a single laboratory. If the results hold, the cost of producing
imitation-learning data falls without the cost being transferred to policy
robustness, which is the trade that passive capture has historically demanded.
